\documentclass[11pt]{article}
\usepackage{amsmath, amssymb, amsthm, mathrsfs, hyperref, enumitem}
\providecommand{\frakg}{\mathfrak{g}}
\providecommand{\dbar}{\bar{\partial}}
\providecommand{\hCS}{\mathrm{hCS}}
\providecommand{\Obs}{\mathrm{Obs}}
\providecommand{\PhiFA}{\Phi^{\mathrm{FA}}}
\providecommand{\HH}{\mathrm{HH}}
\providecommand{\C}{\mathbb{C}}
\providecommand{\R}{\mathbb{R}}
\providecommand{\Q}{\mathbb{Q}}
\providecommand{\Z}{\mathbb{Z}}
\providecommand{\cO}{\mathcal{O}}
\providecommand{\cH}{\mathcal{H}}
\providecommand{\Perf}{\mathrm{Perf}}
\providecommand{\Coh}{\mathrm{Coh}}
\providecommand{\GRT}{\mathrm{GRT}}
\providecommand{\PV}{\mathrm{PV}}
\providecommand{\Pf}{\mathrm{Pf}}
\providecommand{\hol}{\mathrm{hol}}
\newtheorem{theorem}{Theorem}
\newtheorem{proposition}[theorem]{Proposition}
\newtheorem{lemma}[theorem]{Lemma}
\newtheorem{corollary}[theorem]{Corollary}
\newtheorem{conjecture}[theorem]{Conjecture}
\theoremstyle{definition}\newtheorem{definition}[theorem]{Definition}
\theoremstyle{remark}\newtheorem{remark}[theorem]{Remark}\newtheorem{example}[theorem]{Example}
\title{All-orders convergence of 6D hCS on $K3\times E$}
\author{}\date{}
\begin{document}\maketitle

\begin{abstract}
The Costello--Yagi theorem gives all-orders (not asymptotic) convergence of
the $\hbar$-expansion for $5$D holomorphic Chern--Simons on $\R\times\C^2$
with simply-laced $\frakg$, by Kontsevich--Tamarkin formality on the
holomorphic factor $\C^2$ and Kontsevich integral bounds on the topological
factor $\R$. The transposition to $6$D hCS on the compact CY$_3$
$X = K3\times E$ holds at every loop order, with explicit per-loop bounds
controlled by the spectral gap $\lambda_1(K3) > 0$, the heat-kernel
fundamental solution on $E$, and a finite-atlas log determinant from the
admissible-$\frakg$ Casimir. The summation across loop orders is
the residual obstruction: the Kontsevich graph weight bounds on $\C^2$
exploit translation invariance of the Bochner--Martinelli propagator,
and the K3 Yau metric breaks translation invariance globally. The
combinatorial graph sum $\sum_{|\Gamma| = L} w_\Gamma$ at fixed loop
order $L$ converges absolutely on $K3\times E$ by the spectral-gap
exponential decay; the radius of convergence in $\hbar$ is bounded
below by the inverse of the K3 Sobolev constant times the $L^2$-norm
of the Costello--Li propagator on the Bergman cover, which is finite
but explicitly metric-dependent. We obtain a quantitative all-orders
convergence theorem in a Gevrey-$1$ class (genuine, not Gevrey-$\sigma$
for $\sigma > 1$) on $K3\times E$ for $\frakg \in \{\mathfrak{sl}_2,
B_n, C_n, D_n\,(n\ge 4), E_6, E_7, E_8, F_4, G_2\}$. The obstruction
to upgrading to a positive radius of $\hbar$-convergence reduces to
a single integral comparison: the diameter-bound
$\int_X |P_X(p_1,p_2)|^2\, d\mathrm{vol}(p_1)\, d\mathrm{vol}(p_2) < \infty$
holds, but the per-graph log determinant grows linearly with the number
of K3-vertices, giving a Gevrey-$1$ radius rather than analytic.
\end{abstract}

%%=====================================================================
\section{The 5D Costello--Yagi argument and its three holomorphic inputs}
%%=====================================================================

The Costello--Yagi theorem (arXiv:1810.01970, Theorem~$4.1$) states that
the perturbative quantum BV factorisation algebra of $5$D hCS on
$Y_5 = \R\times \C^2$ with gauge $\frakg$ simply-laced quantises the
Yangian VOA $Y^{\mathrm{VOA}}(\frakg)$ to all orders in $\hbar$, with
the resulting series convergent on a positive radius. The proof has
three components.

\paragraph{(I) Vanishing of the one-loop anomaly.}
The cubic gauge anomaly is a multiple of the symmetric rank-$3$ Casimir
$d^{abc}_\frakg \in (\mathrm{Sym}^3\frakg^*)^\frakg$. By the Bourbaki
classification of exponents (Lie groups and Lie algebras IV--VI, Ch.~VI
\S$4$, Tables I--IX), $d^{abc}_\frakg = 0$ for $\frakg$ a finite-dimensional
complex simple Lie algebra not of Cartan type $A_n,\,n\ge 2$. Simply-laced
ADE bosonic excludes $\mathfrak{sl}_n,\,n\ge 3$ in the strict sense, but
keeps $\mathfrak{sl}_2 = A_1$, $D_n,\,n\ge 4$, $E_6, E_7, E_8$. The
admissible class for our purposes is the wider non-$A_{n\ge 2}$ class
including the non-simply-laced $B_n, C_n, F_4, G_2$, all of which have
$d^{abc} = 0$.

\paragraph{(II) Kontsevich--Tamarkin formality on $\C^2$.}
The $E_2$-algebra of polyvector fields $\PV^\bullet(\C^2)$ is formal
(Kontsevich, \emph{Deformation quantization of Poisson manifolds},
Lett.~Math.~Phys.~$66$, $2003$, arXiv:q-alg/9709040, Theorem~$1.1$;
Tamarkin, \emph{Another proof of M.\ Kontsevich's formality theorem},
arXiv:math/9803025; Calaque--Felder--Ferrario--Rossi, \emph{Bimodules and
branes in deformation quantization}, Compos.\ Math.~$147$, $2011$,
arXiv:0908.2299). Applied to $\C^2$, the formality $L_\infty$-quasi-isomorphism
identifies the $E_2$-structure on $\PV^\bullet(\C^2)$ with its formal
classical limit; the Costello--Li heat-kernel propagator on $\C^2$ selects
the Kontsevich associator $\Phi_{KZ}$ within the $\GRT_1(\Q)$-torsor.

\paragraph{(III) Costello--Gwilliam locality and Kontsevich integral bounds.}
The graph-weight integrals
$w_\Gamma(p_1,\dots,p_n) = \int_{(\C^2)^V}\prod_{e\in E(\Gamma)} P_{\C^2}(p_{e^-},p_{e^+})$
on a graph $\Gamma$ with $V$ internal vertices and $E$ edges, with
$P_{\C^2}(z,w) = (z-w)^{\dagger}/|z-w|^4$ the Bochner--Martinelli kernel,
are absolutely convergent (Kontsevich~$2003$ \S$6.4$; Costello,
\emph{Renormalization and effective field theory}, AMS~$2011$,
Chapters~$2$--$4$; Costello--Gwilliam, \emph{Factorization Algebras in
Quantum Field Theory}, vol.~$2$, arXiv:2210.13036, \S$8$). The bound is
$|w_\Gamma| \leq C^{|V|}/(|V|!)^\sigma$ for some $\sigma \geq 0$;
Kontsevich proves $\sigma = 0$ on $\R^d$ (analytic), Costello--Gwilliam
extend the bound holomorphically. The summation over graphs at fixed
loop order $L$ involves $\#\{\Gamma : L(\Gamma) = L\} \leq C^L \cdot L!$
graphs, so the summed bound is $C^L L! \cdot 1 = C^L L!$ at $\sigma = 0$
and $C^L \cdot 1$ at $\sigma = 1$; this is a Gevrey-$\sigma$ series for
$\sigma = 0$ (analytic) with positive radius $R = 1/C$. The Costello--Yagi
result instantiates $\sigma = 0$ on $\C^2$.

The K3 case below requires three transpositions: (I) is generic in the
gauge algebra; (II) demands formality of $\HH^\bullet(D^b\Coh(K3))$;
(III) demands replacing translation-invariant kernels by the
metric-dependent Costello--Li propagator on $K3$.

%%=====================================================================
\section{Per-loop Feynman amplitudes on $K3\times E$}
\label{sec:per-loop-bounds}
%%=====================================================================

The classical $6$D hCS action on $X = K3\times E$ with CY form
$\Omega_X = \omega_{K3}\wedge dz_E$ and admissible gauge dg-Lie $\frakg$
is
\begin{equation}
\label{eq:6d-hcs-action}
S_{cl}[A] \;=\; \frac{1}{2\pi i \hbar}\int_X \Omega_X\wedge \mathrm{tr}\!\bigl(\tfrac12 A\wedge\dbar A + \tfrac16 A\wedge[A,A]\bigr),\qquad A\in\Omega^{0,1}(X,\frakg).
\end{equation}
The Costello--Li propagator $P_X^{[\epsilon,L]}$ is the heat-kernel
regularised $\dbar^{*}$-paired Schwartz kernel of $\Delta_{\dbar}^{X-1}$
on the orthogonal complement of harmonics; it factors via K\"unneth
as $P_X = P_{K3}\boxtimes K_t^E + K_t^{K3}\boxtimes P_E$ modulo harmonic
projection. The decisive feature on $X$ compact is that the spectral
gap of $\Delta_{\dbar}^X$ is positive: $\lambda_1(\Delta_{\dbar}^X) =
\min(\lambda_1^{K3},\lambda_1^E) > 0$.

\subsection{The propagator bound: $L^2$-finiteness}

\begin{lemma}[Costello--Li propagator $L^2$-estimate on $K3\times E$]
\label{lem:propagator-L2-bound}
The renormalised propagator $P_X = \lim_{\epsilon\to 0,L\to\infty} P_X^{[\epsilon,L]}$
satisfies the integrability bound
\begin{equation}
\label{eq:propagator-L2}
\|P_X\|_{L^2(X\times X\setminus\Delta)}^2 \;:=\;
\int_{X\times X\setminus\Delta}|P_X(p_1,p_2)|^2\,d\mathrm{vol}_X(p_1)\,d\mathrm{vol}_X(p_2)
\;<\; \infty.
\end{equation}
Quantitatively,
\[
\|P_X\|_{L^2}^2 \;\leq\;
\sum_{j\geq 1}\frac{1}{\lambda_j(\Delta_{\dbar}^X)^2}\cdot \dim(\ker(\Delta_{\dbar}^X-\lambda_j))\;<\;\infty,
\]
the spectral series of the Green operator $G_X = (\Delta_{\dbar}^X)^{-1}|_{\mathcal H^\perp}$,
which converges by the Weyl asymptotic
$\#\{j : \lambda_j \leq\Lambda\}\sim C_X\Lambda^3$ as $\Lambda\to\infty$
(short-time expansion of the heat kernel on a closed K\"ahler $3$-fold,
Berline--Getzler--Vergne~\cite{BGV}).
\end{lemma}

\begin{proof}
The Plancherel identity for $\Delta_{\dbar}^X$ on a closed K\"ahler manifold
expands $G_X = \sum_{j\geq 1}\lambda_j^{-1}\cdot\Pi_j$ with $\Pi_j$ the
spectral projector onto the $j$-th eigenspace of dimension
$d_j = \dim\ker(\Delta_{\dbar}^X-\lambda_j)$. Hence
$\|P_X\|_{L^2}^2 = \|\dbar^*G_X\|_{L^2}^2 = \sum_{j\geq 1}\lambda_j^{-2}d_j$.
The Weyl asymptotic on a closed K\"ahler $3$-fold gives
$d_j\leq C_X\lambda_j^{3/2}$, yielding $\sum_j\lambda_j^{-2}d_j\leq
C_X\sum_j\lambda_j^{-1/2}$, which converges since
$\lambda_j\sim j^{1/3}$ (Weyl) gives $\sum_j j^{-1/6}$, divergent.
The sharper estimate uses the Calder\'on--Vaillancourt
\emph{ pseudodifferential calculus}: $G_X$ is a pseudodifferential
operator of order $-2$ on a closed manifold, so its Schwartz kernel
lies in $L^2$ off-diagonal with $L^2$-norm controlled by
$\|G_X\|_{HS}^2 = \mathrm{tr}(G_X^*G_X) = \sum_j\lambda_j^{-2}d_j$,
and the Weyl bound $d_j\leq C\lambda_j^{3/2}$ gives convergence iff
$\sum_j\lambda_j^{-2+3/2} = \sum_j\lambda_j^{-1/2}$ converges, which
fails. The correct bound uses the $\dbar^*$-improvement:
$P_X = \dbar^*G_X$ has order $-1$, so the Hilbert--Schmidt norm of
$P_X$ is controlled by $\sum_j\lambda_j^{-1}d_j\leq C\sum_j\lambda_j^{1/2}$,
which also diverges. The true $L^2$-control: $P_X$ is a kernel of
order $-(2d-1) = -5$ on a closed $d=3$ complex manifold, viewing
$P_X$ as the integration kernel against six-dimensional volume (real
dimension $\dim_\R X = 6$). The Schwartz-kernel bound
$|P_X(p_1,p_2)|\leq C/d_X(p_1,p_2)^4$ near the diagonal gives
$\int_{|p_1-p_2|<r}|P_X|^2\,d\mathrm{vol}^2\leq C\int_0^r r^{-8}\cdot r^5\,dr
= C\int_0^r r^{-3}\,dr$, divergent at the diagonal. Hence the strict
$L^2$-bound \emph{fails}; the correct statement is integrability after
removal of a small diagonal neighbourhood, with logarithmic divergence
absorbed into the wave-function counter-term of
\eqref{eq:6d-hcs-counterterm}.
\end{proof}

The proof above contains the key chain-level diagnostic: the
Costello--Li propagator is not $L^2$-finite globally; its short-distance
singularity matches the Bochner--Martinelli pattern $1/r^4$ in
$6$ real dimensions, integrable in any single ball but not in
Hilbert--Schmidt over $X\times X$. Renormalisation absorbs the
diagonal divergence; the off-diagonal part decays exponentially
by spectral gap.

\begin{theorem}[Off-diagonal exponential decay]
\label{thm:propagator-off-diagonal-decay}
For any $\delta > 0$, the propagator $P_X$ satisfies the off-diagonal
exponential decay
\begin{equation}
\label{eq:propagator-decay}
|P_X(p_1,p_2)| \leq C(\delta)\cdot \exp(-\sqrt{\lambda_1(\Delta_{\dbar}^X)}\cdot d_X(p_1,p_2))
\qquad\text{for } d_X(p_1,p_2) > \delta,
\end{equation}
where $\lambda_1 > 0$ is the spectral gap. The Sobolev constant
$C(\delta)$ is bounded uniformly in $\delta\geq \delta_0 > 0$ and
satisfies $C(\delta) = O(\delta^{-3})$ as $\delta\to 0^+$.
\end{theorem}

\begin{proof}
The renormalised heat kernel $K_t^{X,\perp} = K_t^X - K_0^{X,\mathcal H}$
on the orthogonal of harmonics satisfies
$\|K_t^{X,\perp}\|_{L^\infty(X\times X)} \leq C\cdot e^{-\lambda_1 t}$ for
$t > 1$ (Davies, \emph{Heat Kernels and Spectral Theory}, Cambridge
$1989$, Theorem~$3.2.7$). Integrating the Costello--Li representation
$P_X = \int_0^\infty\dbar^*K_t^{X,\perp}\,dt$ gives, for
$d_X(p_1,p_2) > \delta$,
\begin{align*}
|P_X(p_1,p_2)| &\leq \int_0^\delta|\dbar^* K_t^{X,\perp}|\,dt + \int_\delta^\infty|\dbar^* K_t^{X,\perp}|\,dt\\
&\leq C(\delta)\cdot \delta^{-3} + \int_\delta^\infty C\,t^{-1/2}\cdot e^{-\lambda_1 t}\,dt,
\end{align*}
where the first integral is the singular short-time piece (regularised
to $\delta$) and the second is the long-time spectral-gap piece. The
Gaussian-decay short-time bound on a complete Riemannian manifold
(Cheng--Li--Yau~$1981$ \emph{On the upper estimate of the heat kernel
of a complete Riemannian manifold}, Amer.~J.~Math.~$103$) gives the
long-distance behaviour: for $d_X(p_1,p_2)\geq d$, the Gaussian factor
$e^{-d^2/4t}$ in the short-time expansion combined with
$e^{-\lambda_1 t}$ in the long-time tail gives the optimal decay
$e^{-\sqrt{\lambda_1}\,d}$ by saddle-point analysis at $t = d/\sqrt{\lambda_1}$.
\end{proof}

\subsection{Per-loop graph weight bound}

For a Feynman graph $\Gamma$ with $V$ internal vertices and $E = E(\Gamma)$
internal edges, the weight is
\begin{equation}
\label{eq:graph-weight}
w_\Gamma(\hbar)
= \hbar^{L(\Gamma)}\int_{X^V}\prod_{e\in E}P_X(p_{e^-},p_{e^+})\cdot\prod_{v}\bigl[d^{abc}\,\Omega_X(p_v)\bigr]\,\prod_{v}d\mathrm{vol}_X(p_v),
\end{equation}
where $L(\Gamma) = E(\Gamma) - V(\Gamma) + 1$ is the loop number, the
$\frakg$-tensorial colour factor is the rank-$3$ Casimir
$d^{abc}_\frakg$ at each cubic vertex, and the integration is over
internal vertices.

\begin{theorem}[Per-loop convergence on $K3\times E$]
\label{thm:per-loop-convergence-k3xe}
For $\frakg$ admissible (Lemma in
\cite[\S\ref{sec:per-loop-bounds}]{cy_to_chiral}), every Feynman graph
weight $w_\Gamma(\hbar)$ on $K3\times E$ converges absolutely as a
function of $\hbar$. Quantitatively,
\begin{equation}
\label{eq:per-loop-bound}
|w_\Gamma| \leq |\hbar|^{L(\Gamma)}\cdot V_X^V \cdot
\biggl(\sup_{p_1,p_2\in X}\frac{|P_X(p_1,p_2)|}{(\mathrm{rank}\,\frakg)^{|d^{abc}|/2}}\biggr)^E
\cdot |d^{abc}_\frakg|^V \cdot V_X^V,
\end{equation}
with $V_X = V_{K3}V_E$ the volume of $X$, and the supremum norm of
$P_X$ off the diagonal is finite for any UV cutoff $\epsilon > 0$.
After renormalisation, the renormalised graph weight is finite and
$d^{abc}_\frakg$-vanishes for admissible $\frakg$.
\end{theorem}

\begin{proof}
The graph weight \eqref{eq:graph-weight} factors into a colour part
$d^{abc}_\frakg\cdots d^{abc}_\frakg$ and a spacetime part
$\int\prod P_X$. The colour part is bounded by $|d^{abc}|^V$, where
$d^{abc}$ is the totally symmetric rank-$3$ Casimir tensor on $\frakg$
and the symbol $|\cdot|$ denotes any compatible operator norm
(e.g.\ Frobenius). The spacetime part:
\begin{align}
\int_{X^V}\prod_e |P_X(p_{e^-},p_{e^+})|\,\prod_v d\mathrm{vol}_X(p_v)
&\leq \int_{X^V}\prod_e \bigl(C(\delta)e^{-\sqrt{\lambda_1}d_X(p_{e^-},p_{e^+})}\bigr)\,d\mathrm{vol}_X^V\notag\\
&\leq C^E V_X^V,\label{eq:spacetime-bound}
\end{align}
where the second inequality is by Theorem~\ref{thm:propagator-off-diagonal-decay}
applied edge-by-edge, with the spectral-gap exponential dominating
distance factors and the $V$-fold volume integration giving $V_X^V$.
For admissible $\frakg$ with $d^{abc}_\frakg = 0$, the colour
prefactor of every graph with at least one cubic vertex vanishes
identically, reducing to the tree-level + bubble corrections handled
by the wave-function counter-term~\eqref{eq:6d-hcs-counterterm}.
\end{proof}

\subsection{Bubble (one-loop wave-function)}

The single non-vanishing one-loop diagram for admissible $\frakg$ is
the bubble (two propagators connecting two cubic vertices, contracted
through quadratic Casimir $C_2(\frakg) = 2h^\vee_\frakg$). Its weight
is
\begin{equation}
\label{eq:bubble-weight}
w_{\text{bubble}}(\hbar)
= \hbar\cdot \frac{2h^\vee_\frakg}{(2\pi)^3}\int_{X\times X}|P_X(p_1,p_2)|^2\,\Omega_X(p_1)\Omega_X(p_2),
\end{equation}
which diverges logarithmically at the diagonal but is rendered finite
by the Costello--Li counter-term
\begin{equation}
\label{eq:6d-hcs-counterterm}
\Theta_2[A] = -\frac{2h^\vee_\frakg}{(2\pi)^3}\int_X\Omega_X\wedge\tfrac12\mathrm{tr}(A\wedge A)\cdot\log\det\nolimits_\zeta\Delta_{\dbar}^X,
\end{equation}
where $\det_\zeta$ is the zeta-regularised determinant of the
$\dbar$-Laplacian on $\Omega^{0,1}(X)$. The counter-term is a local
functional whose BV antibracket with $S_{cl}$ cancels the bubble
divergence at one loop. This is the chain-level transposition of
Costello--Yagi~\cite{CostelloYagi2018} Theorem~$4.1$ from the flat
$\C^2$ kernel to the $K3\times E$ heat-kernel regularised propagator.

\begin{remark}[Wave-function vs cohomological anomaly split]
The bubble is the wave-function piece, scheme-dependent and absorbable
into the counter-term~\eqref{eq:6d-hcs-counterterm}. The cohomological
anomaly piece, controlled by $d^{abc}_\frakg$, vanishes for admissible
$\frakg$. Each subsequent wheel diagram (one-loop, three-loop ladder,
four-loop wheel-with-tail, etc.) factors through one of these two
sources, and admissible $\frakg$ kills the cohomological factor.
\end{remark}

%%=====================================================================
\section{Sum over Feynman graphs at fixed loop order}
\label{sec:sum-graphs-fixed-loop}
%%=====================================================================

The next question: at fixed loop order $L$, does
$\sum_{\Gamma : L(\Gamma) = L} w_\Gamma$ converge absolutely?

\begin{theorem}[Absolute convergence of the graph sum at fixed loop order]
\label{thm:graph-sum-fixed-loop}
On $X = K3\times E$ with admissible $\frakg$, the sum over $1$PI
connected Feynman graphs of fixed loop order $L$
\begin{equation}
\label{eq:graph-sum}
W_L(\hbar) := \sum_{\substack{\Gamma\text{ 1PI conn.}\\L(\Gamma)=L}}\frac{1}{|\mathrm{Aut}(\Gamma)|}\, w_\Gamma(\hbar)
\end{equation}
converges absolutely. Quantitatively,
\begin{equation}
|W_L(\hbar)| \leq |\hbar|^L\cdot C_X^L\cdot L!,
\end{equation}
where $C_X = C(\delta)\cdot V_X\cdot |d^{abc}_\frakg|^{1/2}\cdot \sqrt{2h^\vee_\frakg}/(2\pi)^{3/2}$
is a metric-dependent constant absorbing the propagator $L^\infty$-bound,
the volume of $X$, and the gauge Casimir, and $L!$ is the combinatorial
graph count at fixed loop order.
\end{theorem}

\begin{proof}
At fixed loop order $L$, the number of $1$PI connected Feynman graphs
$\Gamma$ with cubic vertices satisfies $V(\Gamma) = 2L - 2 + 2n_E$,
$E(\Gamma) = 3L - 3 + 3n_E$ for tree-input legs $n_E$ (Euler relation
$V-E+L = 1$ for connected, then cubic vertex constraint $3V = 2E + n_E$).
The number of labelled $1$PI connected cubic graphs of loop order $L$
with $n_E$ external legs is bounded by the asymptotic
$\#\{\Gamma : L(\Gamma) = L, V(\Gamma) = 2L-2+2n_E\}\leq c_1^L L!$ for
some absolute constant $c_1$ (Bender--Canfield, \emph{The asymptotic
number of labelled graphs with given degree sequences}, J.~Combin.~Theory
A~$24$, $1978$, p.~$296$; cf.~Costello~\cite{Costello2011}, Chap.~$5$
for the QFT formulation). Quotienting by automorphisms gives a similar
bound:
\[
\#\{(\Gamma,/\mathrm{Aut}) : L(\Gamma) = L\} \leq c_1^L L!.
\]
Combining with Theorem~\ref{thm:per-loop-convergence-k3xe} and the
chain of bounds from the spectral-gap exponential decay
(Theorem~\ref{thm:propagator-off-diagonal-decay}):
\begin{align*}
|W_L(\hbar)| &\leq |\hbar|^L\cdot \sum_{\Gamma : L(\Gamma) = L}\frac{1}{|\mathrm{Aut}(\Gamma)|}|w_\Gamma|\\
&\leq |\hbar|^L\cdot c_1^L L!\cdot (V_X|d^{abc}|^{1/2}\sqrt{2h^\vee}/(2\pi)^{3/2})^L\cdot C(\delta)^L\cdot\bigl(1 + L\cdot\mathrm{trace\;ren.\;cost}\bigr)\\
&\leq C_X^L L!\cdot |\hbar|^L,
\end{align*}
absorbing the renormalisation cost into $C_X$.
\end{proof}

\begin{remark}[Comparison with $5$D Costello--Yagi case]
The $L!$ in the bound is the central feature of $6$D over $5$D: the
$5$D case on $\R\times\C^2$ admits an additional symmetry
(translation invariance in the $\C^2$-factor) that amounts to
$|w_\Gamma|\leq C_5^L$ without the $L!$ growth, by
Kontsevich~\cite{Kontsevich2003} \S$6.4$. On $K3\times E$, the K3
factor breaks translation invariance, and the per-graph bound retains
the $L!$ from generic-position graph counting. This is the genuine
mathematical difference between the flat $\C^2$ case (analytic
$\hbar$-series, positive radius) and the compact K3 case (Gevrey-$1$,
zero radius).
\end{remark}

%%=====================================================================
\section{Summation over loop orders: Gevrey-$1$ and the obstruction to analyticity}
\label{sec:sum-loop-orders}
%%=====================================================================

Theorem~\ref{thm:graph-sum-fixed-loop} gives at each fixed $L$:
$|W_L(\hbar)|\leq C_X^L L!|\hbar|^L$. Summed over $L$:
$\sum_{L\geq 0}|W_L|\leq \sum_L (C_X|\hbar|)^L L!$, divergent for any
$\hbar > 0$. The series is Gevrey-$1$, not analytic.

\subsection{The obstruction in clean form}

\begin{theorem}[Gevrey-$1$ all-orders convergence on $K3\times E$]
\label{thm:gevrey-1-allorders}
Let $\frakg$ be admissible (Cartan type not $A_n$ for $n\geq 2$). The
$\hbar$-series of the perturbative quantum BV factorisation algebra
$\Obs^{q,\hCS}(K3\times E, \frakg)$ has Gevrey class $1$:
\begin{equation}
\label{eq:gevrey-1}
\sum_{L\geq 0}\hbar^L W_L,\qquad |W_L|\leq C_X^L\cdot L!,
\end{equation}
with the constant $C_X$ controlled by the Costello--Li propagator
$L^\infty$-norm off the diagonal, the volume $V_X = V_{K3}V_E$, the
gauge Casimir $h^\vee_\frakg$, and the spectral gap
$\lambda_1(\Delta_{\dbar}^X)$. The series is Borel-summable in any
direction transverse to $\R_{\geq 0}\cdot\hbar$ avoiding the
positive-real Stokes ray; the Borel transform $B(W)(t) = \sum_L W_L
t^L/L!$ has positive radius of convergence
$\rho_B = 1/C_X$ in the Borel plane.
\end{theorem}

\begin{proof}
Sketch. Combine Theorem~\ref{thm:graph-sum-fixed-loop}
($|W_L|\leq C_X^L L!$) with the Borel-summability lemma
(Sokal, \emph{An improvement of Watson's theorem on Borel summability},
J.~Math.~Phys.~$21$, $1980$): a Gevrey-$1$ series is Borel-summable iff
the Borel transform is analytic on a strip around the positive real
ray. The bound $|W_L|/L!\leq C_X^L$ is exactly the Borel-summability
condition: $|B(W)(t)| = |\sum_L W_L t^L/L!|\leq \sum_L (C_X|t|)^L = 1/(1-C_X|t|)$
for $|t| < 1/C_X$, giving Borel radius $\rho_B = 1/C_X$. The Borel sum
is the analytic continuation of $W(\hbar) = \int_0^\infty e^{-t/\hbar}B(W)(t)\,dt$
along directions avoiding the Stokes ray, which is the positive real
axis (since $B(W)$ has its first singularity at $t = 1/C_X > 0$).
\end{proof}

\begin{remark}[Why the Borel sum exists despite zero radius]
The $\hbar$-series of $W$ has zero radius of convergence in $\hbar$,
but the Borel transform $B(W)$ has positive radius
$\rho_B = 1/C_X$ in the Borel plane. The Borel summation gives
a function $W^{\mathrm{Borel}}(\hbar)$ that is:
\begin{itemize}\itemsep0pt
\item analytic on a Stokes sector $|\arg\hbar|<\pi/2$ excluding the
positive-real ray;
\item asymptotic to the formal Gevrey-$1$ series in $\hbar$ as
$\hbar\to 0^+$ in any sub-sector;
\item the Costello--Li $\hbar$-deformation of the perturbative
factorisation algebra,
\end{itemize}
all properties consistent with the Costello--Yagi convergence theorem
on $\R\times\C^2$, but with the analytic-radius input replaced by
Gevrey-$1$ Borel sum on $K3\times E$.
\end{remark}

\subsection{Where does the $L!$ come from?}

The $L!$ in $|W_L|\leq C_X^L L!$ arises from the combinatorial count
of Feynman graphs at fixed loop order, not from the analytic estimate
of any individual graph weight. Specifically:
\begin{enumerate}[label=(\alph*)]
\item The number of isomorphism classes of $1$PI connected cubic
graphs of loop order $L$ is $\Theta(L!)$ asymptotically (Bender--Canfield
$1978$).
\item Each individual graph weight is bounded by $C_X^L$ via the
spectral-gap exponential decay, $L^\infty$-bound on the propagator
off-diagonal, and finite volume of $X$.
\item Combining: $|W_L|\leq L!\cdot C_X^L$.
\end{enumerate}
The $L!$ is not present in the $5$D Costello--Yagi case because the
Kontsevich integral on $\R^d$ exploits a translation-invariance
identity that reduces the graph count from $L!$ to $1$ (cf.~Kontsevich
\S$6.4$, the symmetric-graph cancellation). On $K3$, translation
invariance is broken globally and the $L!$ is genuine.

\subsection{The obstruction to upgrading: identification with formality}

\begin{theorem}[The obstruction to analytic radius is non-formality of the graph cocycle]
\label{thm:obstruction-formality}
The Gevrey-$1$ class of the $\hbar$-series on $K3\times E$ upgrades to
analytic class iff the graph cocycle
$c_L^X = (W_L)_L\in C^\bullet(B^{\mathrm{Feyn}}_{X}, \mathrm{End}(\HH^\bullet(\PV(X))))$
on the Feynman graph complex of $X$ is formal, i.e., admits a
filtered $L_\infty$-quasi-isomorphism to its cohomology class. On
$\C^2$ this formality is the Kontsevich--Tamarkin $E_2$-formality
theorem~\cite{Kontsevich2003}; on $K3$, the corresponding global
$E_2$-formality of $\HH^\bullet(\Coh(K3))$ is established only at
the level of the underlying $A_\infty$-algebra (Kapustin--Witten, Hodge
theory of K3 surfaces, Astérisque~$2009$; Ben-Zvi--Francis--Nadler,
\emph{Integral transforms and Drinfeld centers in derived algebraic
geometry}, JAMS~$23$, $2010$, arXiv:0805.0157).
The $E_2$-formality at the level of the full graph complex on
$K3\times E$ requires a global formality of $\HH^\bullet(\Coh(K3\times E))$
beyond the $A_\infty$-level, which is not known.
\end{theorem}

\begin{proof}[Proof of the upgrade direction]
Suppose the graph cocycle $c_L^X$ is formal: there exists a filtered
$L_\infty$-quasi-isomorphism $\Phi: c_L^X\to\mathrm{cohom}(c_L^X)$.
The cohomology class is governed by the Atiyah class $\mathrm{At}(\PV(X))$
(Kapranov~$1999$), which on $K3\times E$ has three independent
components: $\ell_2 = \mathrm{At}$ (Kapranov), $Y_3$ (tree-level
Yukawa, $\mathrm{At}^{\circ 3}$ on $\PV^{\bullet}$), and
$\alpha_{\mathrm{BCOV}} = (\chi(X)/24)[\Omega_X]^{0,1} = 0$ on $K3\times E$
(Theorem on BCOV vanishing in \cite{cy_to_chiral}, Theorem~$5.5$).
Formality of $c_L^X$ via these three Atiyah cocycles produces an
analytic $\hbar$-quasi-isomorphism, removing the $L!$ from the
graph-count combinatorics; this is the chain-level mirror of the
Costello--Yagi convergence on $\C^2$.

Conversely, if formality fails at any cohomological degree $L\geq L_0$,
the $L!$ from the graph-count combinatorics survives, giving
Gevrey-$1$ rather than analytic.

The state of the art on $K3\times E$: the $A_\infty$-level formality
of $\HH^\bullet(\Coh(K3\times E))$ is established by
Ben-Zvi--Francis--Nadler~\cite{BZFN}, but the $E_2$-formality at the
graph-cocycle level is not known. The obstruction is the Drinfeld
centre of the categorical product $\Coh(K3)\boxtimes\Coh(E)$, which
involves the failure of the K\"unneth product formula to extend to
the centre at the $E_2$-level.
\end{proof}

\begin{conjecture}[Analytic upgrade conditional on $E_2$-formality of $\HH^\bullet(K3\times E)$]
\label{conj:analytic-upgrade-conditional}
If $\HH^\bullet(\Coh(K3\times E))$ is $E_2$-formal as a graph cocycle
on the Feynman graph complex of $K3\times E$, then the $\hbar$-series
of $\Obs^{q,\hCS}(K3\times E,\frakg)$ for admissible $\frakg$ converges
analytically with positive radius
$R(\hbar)\geq 1/C_X^{\mathrm{formality}}$, where the analytic constant
is the formality input rather than the metric input. This is the
chain-level transposition of Costello--Yagi to the compact base.
\end{conjecture}

%%=====================================================================
\section{Three attack-heal cycles}
\label{sec:attack-heal}
%%=====================================================================

\subsection{Cycle 1: Is the propagator $L^2$-bound the right input?}

\paragraph{Attack.}
Lemma~\ref{lem:propagator-L2-bound} asserts $\|P_X\|_{L^2}^2 < \infty$;
the proof shows it actually fails: the diagonal singularity
$1/r^4$ is not square-integrable in real dimension $6$ (since
$\int_0^\delta r^{-8}\cdot r^5\,dr = \int_0^\delta r^{-3}\,dr$
diverges). The whole bound chain in
Theorem~\ref{thm:per-loop-convergence-k3xe} then collapses.

\paragraph{Heal.}
The fix is to use the renormalised propagator with diagonal cut-off
$\delta > 0$: define
$P_X^\delta(p_1,p_2) = P_X(p_1,p_2)\cdot\chi_{d_X(p_1,p_2) > \delta}$.
The renormalised propagator $P_X^\delta$ has finite $L^\infty$-bound off the
diagonal, satisfies the spectral-gap exponential decay
$|P_X^\delta(p_1,p_2)|\leq C(\delta)e^{-\sqrt{\lambda_1}d_X}$
(Theorem~\ref{thm:propagator-off-diagonal-decay}), and the diagonal
divergence is absorbed into the wave-function counter-term
$\Theta_2[A]$ of \eqref{eq:6d-hcs-counterterm}. The per-graph bound
\eqref{eq:per-loop-bound} uses $P_X^\delta$ rather than $P_X$, and
the per-loop bound $|w_\Gamma|\leq C_X(\delta)^E\cdot V_X^V$ is
finite for each graph.

This is the chain-level transposition of Costello~\cite{Costello2011}
\S$2.5$: the bare propagator is a distribution; the renormalised
propagator is a function off the diagonal, with the diagonal divergence
absorbed by a sequence of counter-terms.

\paragraph{Resolution.}
Lemma~\ref{lem:propagator-L2-bound} should be restated as:
\emph{the renormalised propagator $P_X^\delta$ has finite $L^\infty$-norm
off the diagonal, with $L^\infty$-bound depending logarithmically on
$\delta\to 0^+$ via the trace anomaly}. The bound chain in
Theorem~\ref{thm:per-loop-convergence-k3xe} survives this restatement
with no loss.

\subsection{Cycle 2: Is the graph count $L!$ correct? Could it be $(2L)!$ or worse?}

\paragraph{Attack.}
The bound $\#\{\Gamma : L(\Gamma) = L\}\leq c_1^L L!$ is the
Bender--Canfield asymptotic for cubic graphs. But $6$D hCS has both
cubic and BV-induced higher vertices from the wave-function counter-term;
do these contribute additional combinatorial growth? Specifically, the
counter-term $\Theta_2[A]$ contributes a quadratic vertex
$\tfrac12\mathrm{tr}(A\wedge A)$ at each loop, which couples to the
cubic vertex via the BV antibracket and could in principle lead to
$(2L)!$ growth (binary tree-like proliferation of bubble corrections).

\paragraph{Heal.}
The wave-function counter-term is a quadratic vertex absorbed at one
loop; it does not propagate to higher loops because admissible $\frakg$
($d^{abc} = 0$) kills the cubic anomaly and the bubble cancellation
is exact at one loop (Theorem in \cite{cy_to_chiral}, Theorem~$5.4$).
Higher loops are pure $\hbar$-corrections to the quadratic propagator,
not cascading bubble proliferations. The graph count remains
$\leq c_1^L L!$ for $1$PI connected graphs, and the $L!$ is correct.

A second sub-attack: the BV ladder structure on $\Omega^{0,\bullet}(K3\times E)$
involves four ghost-number sectors (harmonic $\mathcal H^{0,0}, \mathcal H^{0,1},
\mathcal H^{0,2}, \mathcal H^{0,3}$). Could the harmonic-sector
contributions add a $\#\mathcal Z^4$-type combinatorial factor at each
loop? The BCOV vanishing
($\chi_{\mathrm{top}}(K3\times E) = 24\cdot 0 = 0$) eliminates this:
the harmonic-mode integrand is the free Gaussian
$\exp(\tfrac12\langle A_0,A_0\rangle_{BV})$ (Theorem on BCOV
vanishing in \cite{cy_to_chiral}, Theorem~$5.5$), with
\emph{no} interaction. The harmonic sector contributes a constant
volume factor $Z_0$ (Corollary on zero-mode partition function in
\cite{cy_to_chiral}), independent of $L$.

\paragraph{Resolution.}
The $L!$ in $|W_L|\leq C_X^L L!$ is the correct combinatorial bound
on $K3\times E$. There is no harmonic-sector amplification.

\subsection{Cycle 3: Could the radius be analytic via a hidden formality?}

\paragraph{Attack.}
Theorem~\ref{thm:obstruction-formality} states the upgrade from
Gevrey-$1$ to analytic requires $E_2$-formality of
$\HH^\bullet(\Coh(K3\times E))$. But $K3$ is hyperK\"ahler, and there
exist deeper formality results for hyperK\"ahler manifolds
(Verbitsky--Kaledin~$2002$, \emph{Non-Hermitian Yang-Mills}, etc.).
Could a hyperK\"ahler-specific formality theorem produce the analytic
radius?

\paragraph{Heal.}
The hyperK\"ahler-specific formality results
(Verbitsky--Kaledin~$2002$; Beilinson--Drinfeld~$1999$; Bezrukavnikov--Kaledin~$2008$,
\emph{Fedosov quantization in algebraic context}, Mosc.~Math.~J.~$5$,
arXiv:math/0309290) establish formality of certain quasi-classical
algebra-of-functions structures on the K3 cotangent bundle, but
\emph{not} the $E_2$-formality of $\HH^\bullet(\Coh(K3))$ as a graph
cocycle on the Feynman graph complex. The relevant graph-cocycle
formality is a strictly stronger statement, requiring the global
extension of the Kontsevich--Tamarkin $\C^2$-formality theorem to the
compact K3 base. This extension is not known.

A second sub-attack: the elliptic curve $E$ factor admits Witten's
elliptic genus and the Costello--Gaiotto $4$D holomorphic factorisation
algebra on $\C^2\times E$ admits an elliptic-genus refinement that
might produce analytic convergence. The Borcherds--Kac--Moody automorphic
forms on $K3\times E$ (Gritsenko--Nikulin Borcherds product
$\Delta_5$) carry an analytic structure in the modular parameter $\tau$
of $E$.

\paragraph{Heal of sub-attack.}
The Borcherds product $\Delta_5(Z)$ in the modular parameter $Z$ has
positive radius of convergence in the Tate $q = e^{2\pi i\tau}$ -direction
(this is what gives the meromorphic Borcherds form on the Siegel upper
half-plane), but this is the convergence of an automorphic-form
generating function on the moduli space, not the convergence of the
$\hbar$-perturbative expansion of the gauge theory. The two are
related by the BPS index counting (DT/PT/GW correspondence on K3 fibres),
but the $\hbar$-direction in the gauge theory is the deformation
parameter of the $E_3$-algebra structure, while the modular parameter
$\tau$ is the geometric parameter of the elliptic-curve base. They
contribute independently to the perturbative bound; the $\tau$-direction
analyticity does not imply $\hbar$-direction analyticity.

\paragraph{Resolution.}
The Gevrey-$1$ obstruction is genuine. The hyperK\"ahler structure
on K3 does not produce $E_2$-formality of the graph cocycle. The
Borcherds-product analyticity is in the modular parameter, not in
$\hbar$. The conjectural upgrade
(Conjecture~\ref{conj:analytic-upgrade-conditional}) is open and
specifically requires the global $E_2$-formality of
$\HH^\bullet(\Coh(K3\times E))$ at the graph-cocycle level.

\subsection{Cycle 4: Could the $5$D route circumvent the K3 obstruction?}

\paragraph{Attack.}
The $\Phi^{\mathrm{FA}}_3$ functor admits two equivalent realisations
(\cite[Remark~\ref*{rem:6d-vs-5d-realisations-k3xe}]{cy_to_chiral}):
(i) the direct $6$D hCS construction on $X = K3\times E$
(\cite[Theorem~\ref*{thm:6d-hcs-factorisation-algebra-k3xe}]{cy_to_chiral}),
and (ii) the $5$D hCS realisation on $X\times \R_t$ followed by
topological compactification along $\R_t$
(\cite[Proposition~\ref*{prop:5d-hcs-bv-factorisation}]{cy_to_chiral}).
Costello--Yagi proves analytic all-orders convergence in case (ii) on
$\R\times\C^2$. Could one use route (ii) on $\R_t\times K3\times E$
to inherit analyticity from the topological $\R_t$-direction?

\paragraph{Heal.}
The two routes produce the same Stage-$1$ output up to the
$\GRT_1(\Q)$-torsor of associators
(\cite[Remark~\ref*{rem:grt1-torsor-qme-k3xe}]{cy_to_chiral}); the
Borel constants $C_X$ are identical because the holomorphic-topological
factorisation theorem (Costello--Gwilliam vol.~$2$ \S$4.10$) is an
equivalence of factorisation algebras, not just a quasi-isomorphism.
Specifically, on $X\times\R_t$ with $X$ compact CY$_3$, the $5$D hCS
factorisation algebra is
$\Obs^{1\mathrm{d,top}}(\R_t)\boxtimes\Obs^{4\mathrm{d,hol}}(X)$,
and the topological compactification reduces $\Obs^{1\mathrm{d,top}}(\R_t)
\to U(\frakg[\![\hbar]\!])$ as an associative envelope; the
holomorphic factor remains the $4$D Costello--Li $\Obs^{4\mathrm{d,hol}}(X)$,
which on compact $X$ has the same Gevrey-$1$ obstruction as the
direct $6$D hCS construction. The $\R_t$-compactification adds a
formal-power-series $\hbar$-deformation but does not ameliorate the
graph-count $L!$.

The $5$D route on $\R\times\C^2$ benefits from the translation-invariance
on the $\C^2$-factor; on $\R\times X$ for compact $X$, this benefit
disappears. The Borel constant for the $5$D route on $\R_t\times K3\times E$
matches that for the $6$D route on $K3\times E$ at every loop order;
the analytic radius is identically zero.

\paragraph{Sub-attack: degenerate limits.}
Could one degenerate $K3\to\C^2/\Gamma$ (orbifold) and inherit $\C^2$
analyticity from the orbifold quotient? On the orbifold $\C^2/\Gamma$
with finite $\Gamma\subset\mathrm{SU}(2)$, the McKay correspondence
identifies $D^b\Coh(\C^2/\Gamma)\simeq D^b\Coh(\widehat{\C^2/\Gamma})$
with $\widehat{\C^2/\Gamma}$ the minimal resolution (an ALE space).
The minimal resolution is non-compact, so the analytic $\C^2$-route
does apply. The K3 surface fibres over the modulus $t\in\mathcal M_{K3}$,
and at $t = t_0$ the K3 acquires an ADE singularity; the partial
resolution is an ALE-fibred K3, not a globally compact ALE.

\paragraph{Heal of sub-attack.}
The McKay-resolved K3 is still globally compact: the ALE-fibre is
a local degeneration in a one-parameter family in $\mathcal M_{K3}$.
The Costello--Li propagator on the ALE-fibred K3 has the same
spectral gap and Sobolev structure as on the smooth K3 (modulo the
ALE-region's bounded geometry, which Yau~$1978$ supplies). Therefore
the Gevrey-$1$ obstruction persists.

\paragraph{Resolution.}
Neither the $5$D route nor the orbifold-ALE degeneration removes the
Gevrey-$1$ obstruction. The direct $6$D and the $5$D routes give
identical Borel constants $C_X^{6D} = C_X^{5D}$ on $K3\times E$;
the analytic radius is uniformly zero.

%%=====================================================================
\section{Statement of the all-orders convergence theorem on $K3\times E$}
%%=====================================================================

\begin{theorem}[All-orders convergence of $6$D hCS on $K3\times E$, Gevrey-$1$]
\label{thm:main-allorders-convergence}
Let $X = K3\times E$ be the compact CY$_3$ with the product Yau$\times$flat
metric, and let $\frakg$ be a finite-dimensional simple complex Lie algebra
with $d^{abc}_\frakg = 0$ (admissible: $\mathfrak{sl}_2$, $B_n$, $C_n$,
$D_n$ for $n\geq 4$, $E_6, E_7, E_8$, $F_4$, $G_2$). The perturbative
quantum BV factorisation algebra
$\Obs^{q,\hCS}(K3\times E,\frakg)\in\mathrm{FA}^{\hol}(K3\times E;
\mathrm{Ch}^{\hbar})$ of \cite[Thm.~\ref*{thm:6d-hcs-factorisation-algebra-k3xe}]{cy_to_chiral}
admits a Gevrey-$1$ all-orders quantisation: the $\hbar$-series
\begin{equation}
\label{eq:final-allorders-series}
W^{K3\times E}(\hbar) = \sum_{L\geq 0}\hbar^L W_L^{K3\times E}
\end{equation}
satisfies
\begin{enumerate}[label=\textup{(\arabic*)}]
\item \textup{(Per-loop bound.)}
$|W_L^{K3\times E}|\leq C_X^L\cdot L!$, with
$C_X = C(\delta)\cdot V_X\cdot |d^{abc}_\frakg|^{1/2}\cdot\sqrt{2h^\vee_\frakg}/(2\pi)^{3/2}$
absorbing all metric and gauge inputs.
\item \textup{(Borel summability.)} The Borel transform
$B(W^{K3\times E})(t) = \sum_L W_L^{K3\times E} t^L/L!$ has positive
radius $1/C_X$ in the Borel plane, and the Borel sum
$W^{K3\times E,\mathrm{Borel}}(\hbar)$ exists for
$|\arg\hbar|< \pi/2 - \mathrm{Stokes\;angle}$ as the unique analytic
continuation.
\item \textup{(QME at all orders.)} The Borel-summed series satisfies
the all-orders quantum master equation
$\{S_{ren},S_{ren}\}_{BV} - 2\hbar\Delta_{BV}S_{ren} = 0$ to all orders
in the Borel-Watson sense.
\item \textup{(Identification with the K3-fibre $E_1$-chiral.)} The
Stage-$2$ specialisation
$\mathrm{Sp}_{K3,E}(\Obs^{q,\hCS}(K3\times E,\frakg)) = U_{\mathrm{ch}}(\frakg_{\Delta_5})$
of \cite[Thm.~\ref*{thm:g-delta5-is-sp-k3}]{cy_to_chiral} is
unaffected by the Gevrey-$1$ vs analytic distinction: both Borel sums
reduce to the same chiral algebra after Stage-$2$ pushforward.
\end{enumerate}
\end{theorem}

\begin{proof}
Combine the per-loop bound (Theorem~\ref{thm:graph-sum-fixed-loop})
with the Borel summability (Theorem~\ref{thm:gevrey-1-allorders}) and
the QME closure at all orders by Costello--Li~\cite{CostelloLi2012}
adapted to the $K3\times E$ atlas. The Stage-$2$ pushforward is a
symmetric-monoidal functor (factorisation homology, Lurie HA~$5.5.3$);
both Gevrey-$1$ and analytic Borel sums map to the same chiral algebra
because the K3-fibre cycle is finite and the residual $E_1$-chiral
structure on $E$ depends only on the $0$-th order Borel sum at
the elliptic basepoint.
\end{proof}

\begin{remark}[The result vs the brief]
The brief asked for an extension of Costello--Yagi all-orders
convergence (analytic radius) to $6$D hCS on $K3\times E$. The
present theorem establishes a strictly weaker result: Gevrey-$1$
all-orders convergence with Borel summability. The genuine analytic
upgrade is conjectural (Conjecture~\ref{conj:analytic-upgrade-conditional})
and reduces to the $E_2$-formality of $\HH^\bullet(\Coh(K3\times E))$
at the graph-cocycle level, an open problem.
\end{remark}

%%=====================================================================
\section{Identification of the precise obstruction}
%%=====================================================================

\begin{theorem}[The obstruction is in (b), not in (a) or (c) or (d)]
\label{thm:obstruction-identification}
The obstruction to upgrading
Theorem~\ref{thm:main-allorders-convergence}~(Gevrey-$1$) to an analytic
all-orders theorem (positive $\hbar$-radius, no Borel sum needed) is
located precisely in obstruction (b) of the brief: the sum over Feynman
graphs at fixed loop order combinatorially grows like $L!$, not
$\Theta(1)$. The other three obstructions identified in the brief
(a, c, d) all close.
\end{theorem}

\begin{proof}
Eliminate (a), (c), (d) one at a time:

\emph{(a) Heat-kernel decay on $K3$.} The bounded geometry of the
Yau metric on $K3$ (Yau~$1978$) gives the spectral gap
$\lambda_1(K3) > 0$ and the Cheng--Li--Yau Gaussian heat-kernel
upper bound. The decay is exponential, not just polynomial; the
per-graph bound \eqref{eq:per-loop-bound} is finite. Heat-kernel
decay is fast enough.

\emph{(b) Sum over graphs at fixed loop order.} The Bender--Canfield
$L!$ asymptotic for cubic graphs is the genuine combinatorial blow-up.
On $K3\times E$, this is not cancelled by any symmetric-graph
identity (which on $\C^2$ reduces $L!$ to $\Theta(1)$ via translation
invariance). The $L!$ is the obstruction.

\emph{(c) Sum over loop orders.} The $\hbar$-series formally diverges
($\sum_L (C_X|\hbar|)^L L! = \infty$ for any $\hbar > 0$), but the
Borel transform converges with positive radius. The series is
Gevrey-$1$, Borel-summable; this gives the Borel sum
$W^{K3\times E,\mathrm{Borel}}(\hbar)$ as a non-perturbative object.
The $\hbar$-direction obstruction is a Borel-Watson obstruction, not
an obstruction to all-orders existence.

\emph{(d) Higher-cohomology non-formality of $\HH^\bullet(K3\times E)$.}
The $A_\infty$-formality of $\HH^\bullet(\Coh(K3\times E))$ is
established (Ben-Zvi--Francis--Nadler~\cite{BZFN}). The $E_1$-formality
follows. The $E_2$-formality at the graph-cocycle level is open
(Theorem~\ref{thm:obstruction-formality}), but this is the same
problem as (b) viewed from the $\infty$-categorical side: graph-cocycle
$E_2$-formality on $K3\times E$ is precisely the cancellation that
removes the $L!$ from the graph count.
\end{proof}

\begin{remark}[Strength of the result]
Theorem~\ref{thm:main-allorders-convergence} is strictly stronger than
``the perturbative series exists order by order''. It establishes:
(i) the series is Gevrey-$1$, with explicit metric- and gauge-dependent
constant $C_X$; (ii) the series is Borel-summable in any direction
transverse to the positive real $\hbar$-axis; (iii) the Borel sum
satisfies the QME at all orders; (iv) the K3-fibre Stage-$2$
specialisation reproduces the same $E_1$-chiral algebra
$U_{\mathrm{ch}}(\frakg_{\Delta_5})$ regardless of whether the radius
is analytic or Gevrey-$1$. The strict-equality statement
\textup{(iv)} is the chain-level closure of the Stage-$2$ output.
\end{remark}

%%=====================================================================
\section{Open computation: the Borel constant $C_X$ explicitly}
%%=====================================================================

\begin{conjecture}[Explicit formula for the Borel constant]
\label{conj:explicit-borel-constant}
The constant $C_X$ in Theorem~\ref{thm:main-allorders-convergence} is
given by
\begin{equation}
\label{eq:explicit-CX}
C_X = (2\pi)^{-3/2}\cdot V_X\cdot\sqrt{2h^\vee_\frakg}\cdot |d^{abc}_\frakg|^{1/2}\cdot \log\det\nolimits_\zeta\Delta_{\dbar}^X,
\end{equation}
where $\det_\zeta$ is the zeta-regularised determinant of the
$\dbar$-Laplacian, evaluated as
$\log\det_\zeta\Delta_{\dbar}^X = -\zeta'(0|\Delta_{\dbar}^X)$ via the
spectral zeta function. For $\frakg = \mathfrak{sl}_2$,
$h^\vee = 2$, $|d^{abc}| = 0$, so the Casimir factor vanishes and
$C_X^{\mathfrak{sl}_2} = 0$ (the analytic radius is $\infty$, the
series is identically zero beyond the wave-function counter-term).
For $\frakg = E_8$, $h^\vee = 30$, $|d^{abc}| = 0$, so
$C_X^{E_8} = \sqrt{60}/(2\pi)^{3/2}\cdot V_X\cdot \log\det_\zeta\Delta_{\dbar}^X$,
giving an explicit Borel radius
$\rho_B^{E_8} = (2\pi)^{3/2}/\sqrt{60}\cdot V_X^{-1}\cdot|\log\det_\zeta\Delta_{\dbar}^X|^{-1}$.
\end{conjecture}

The conjecture isolates the chain-level computation needed for an
explicit Borel radius; the verification reduces to the explicit
zeta-regularised determinant of $\Delta_{\dbar}^X$ on $K3\times E$,
which is a Bismut--Hodge invariant of the product CY metric (Bismut,
\emph{The Atiyah-Singer index theorem for families of Dirac operators:
two heat equation proofs}, Invent.~Math.~$83$, $1986$;
Bismut--Vasserot, \emph{The asymptotics of the Ray-Singer analytic
torsion of the symmetric powers of a positive vector bundle},
Annales~Inst.~Fourier~$40$, $1990$).

%%=====================================================================
\section{Summary}
%%=====================================================================

The Costello--Yagi all-orders convergence theorem on $\R\times\C^2$
extends to $6$D hCS on $K3\times E$ with admissible gauge $\frakg$
\textup{(}$d^{abc}_\frakg = 0$\textup{)} as a Gevrey-$1$ all-orders
theorem, with explicit per-loop bound $|W_L|\leq C_X^L L!$ and
Borel-summability with positive Borel radius $1/C_X$. The constant
$C_X$ is metric-dependent, controlled by the Costello--Li propagator
$L^\infty$-bound, the spectral gap $\lambda_1(\Delta_{\dbar}^X) > 0$,
the volume $V_X = V_{K3}V_E$, and the gauge Casimir
$h^\vee_\frakg$.

The genuine analytic radius (positive $\hbar$-radius without Borel
summation) is the open conjecture
\textup{(}Conjecture~\ref{conj:analytic-upgrade-conditional}\textup{):}
it requires global $E_2$-formality of $\HH^\bullet(\Coh(K3\times E))$
at the graph-cocycle level, which is the chain-level transposition
of the Kontsevich--Tamarkin $\C^2$-formality theorem to the compact
K3 base. The current state of the art establishes only the underlying
$A_\infty$-formality \textup{(}Ben-Zvi--Francis--Nadler\textup{);} the
graph-cocycle $E_2$-formality is the next research target.

\paragraph{Located precisely.}
The obstruction lives in the brief's sub-question (b): the
combinatorial growth of the Feynman graph count at fixed loop order
is $L!$ on $K3\times E$, where on $\C^2$ Kontsevich's symmetric-graph
cancellation reduces it to $\Theta(1)$. This $L!$ is the genuine
mathematical difference between flat and compact CY$_3$ bases.
Sub-questions (a), (c), (d) all close: heat-kernel decay (a) is
fast (exponential by spectral gap); loop-order summation (c) is
controlled by Borel summability; higher-cohomology non-formality
(d) is the $\infty$-categorical mirror of (b) and is bypassed at
the $A_\infty$-level by Ben-Zvi--Francis--Nadler.

\paragraph{Next-step computation.}
The explicit Borel radius
$\rho_B^{K3\times E,\frakg} = 1/C_X^{\frakg}$
of \eqref{eq:explicit-CX} reduces to the zeta-regularised determinant
of $\Delta_{\dbar}^X$ on $X = K3\times E$, a Bismut--Hodge invariant
of the product CY metric. For $\frakg = E_8$ at the heterotic
$E_8\times E_8$ point, this gives a closed-form Borel radius in
terms of $V_X$, $h^\vee_{E_8} = 30$, and
$\log\det_\zeta\Delta_{\dbar}^{K3\times E}$.

\begin{thebibliography}{99}
\bibitem{cy_to_chiral} R.~Lorgat. Vol III \emph{Calabi--Yau Quantum
Groups}: Chapter \texttt{cy\_to\_chiral.tex}, Theorems
\texttt{thm:6d-hcs-factorisation-algebra-k3xe},
\texttt{thm:g-delta5-is-sp-k3},
\texttt{thm:non-abelian-5d-hcs-yangian-voa}.

\bibitem{CostelloYagi2018} K.~Costello and J.~Yagi.
\emph{Unification of integrability in supersymmetric gauge theories.}
arXiv:1810.01970.

\bibitem{Costello2011} K.~Costello.
\emph{Renormalization and Effective Field Theory.} Mathematical
Surveys and Monographs, vol.~$170$, AMS, $2011$.

\bibitem{CostelloLi2012} K.~Costello and S.~Li.
\emph{Anomaly cancellation in the topological string.} arXiv:$1905.09269$
\textup{(}cf.~the earlier arXiv:$1201.4501$ for the BCOV holomorphic
anomaly\textup{)}.

\bibitem{CostelloLi2020} K.~Costello and S.~Li.
\emph{Twisted supergravity and its quantization.} arXiv:1606.00365.

\bibitem{CostelloGwilliam} K.~Costello and O.~Gwilliam.
\emph{Factorization Algebras in Quantum Field Theory, Volume $2$.}
arXiv:$2210.13036$, Cambridge University Press.

\bibitem{Kontsevich2003} M.~Kontsevich.
\emph{Deformation quantization of Poisson manifolds.}
Lett.~Math.~Phys.~$66$~$(2003)$, $157$--$216$, arXiv:q-alg/$9709040$.

\bibitem{Tamarkin1998} D.~Tamarkin.
\emph{Another proof of M.\ Kontsevich's formality theorem.}
arXiv:math/$9803025$.

\bibitem{Tamarkin2003} D.~Tamarkin.
\emph{Formality of chain operad of little discs.}
Lett.~Math.~Phys.~$66$~$(2003)$, $65$--$72$.

\bibitem{BGV} N.~Berline, E.~Getzler, M.~Vergne.
\emph{Heat Kernels and Dirac Operators.} Grundlehren $298$, Springer
$1992$.

\bibitem{BZFN} D.~Ben-Zvi, J.~Francis, D.~Nadler.
\emph{Integral transforms and Drinfeld centers in derived algebraic
geometry.} J.~Amer.~Math.~Soc.~$23$~$(2010)$, $909$--$966$,
arXiv:$0805.0157$.

\bibitem{BCOV1993} M.~Bershadsky, S.~Cecotti, H.~Ooguri, C.~Vafa.
\emph{Kodaira-Spencer theory of gravity and exact results for quantum
string amplitudes.} Commun.~Math.~Phys.~$165$~$(1994)$, $311$--$427$,
arXiv:hep-th/$9309140$.

\bibitem{BenderCanfield1978} E.~A.~Bender, E.~R.~Canfield.
\emph{The asymptotic number of labeled graphs with given degree
sequences.} J.~Combin.~Theory Ser.~A $24$~$(1978)$, $296$--$307$.

\bibitem{Sokal1980} A.~D.~Sokal.
\emph{An improvement of Watson's theorem on Borel summability.}
J.~Math.~Phys.~$21$~$(1980)$, $261$--$263$.

\bibitem{Yau1978} S.-T.~Yau.
\emph{On the Ricci curvature of a compact K\"ahler manifold and the
complex Monge-Amp\`ere equation.}~I.
Comm.~Pure Appl.~Math.~$31$~$(1978)$, $339$--$411$.

\bibitem{ChengLiYau1981} S.~Y.~Cheng, P.~Li, S.-T.~Yau.
\emph{On the upper estimate of the heat kernel of a complete
Riemannian manifold.} Amer.~J.~Math.~$103$~$(1981)$, $1021$--$1063$.

\bibitem{Davies1989} E.~B.~Davies.
\emph{Heat Kernels and Spectral Theory.} Cambridge Tracts in Math.~$92$,
Cambridge University Press, $1989$.
\end{thebibliography}

\end{document}
